% Generated draft: requires mathematical and source review.
\documentclass[11pt]{article}
\usepackage[margin=25mm]{geometry}
\usepackage{fontspec}
\setmainfont{Latin Modern Roman}
\usepackage{amsmath,amssymb,mathtools,unicode-math}
\setmathfont{Latin Modern Math}
\usepackage{enumitem,needspace,xurl,hyperref}
\setlength{\parindent}{0pt}
\setlength{\parskip}{5pt}
\setlength{\emergencystretch}{3em}
\newcommand{\N}{\mathbb{N}}
\newcommand{\Z}{\mathbb{Z}}
\newcommand{\Q}{\mathbb{Q}}
\newcommand{\R}{\mathbb{R}}
\newcommand{\C}{\mathbb{C}}
\newcommand{\F}{\mathbb{F}}
\newcommand{\A}{\mathbb{A}}
\newcommand{\PP}{\mathbb{P}}
\newcommand{\E}{\mathbb{E}}
\newcommand{\eps}{\varepsilon}
\newcommand{\dd}{\,\mathrm{d}}
\begin{document}
\section*{UnsolvedMath ID: 30002367}
\textbf{Source problem code:} catalogue entry
\subsection{Definitions and mathematical statement}
Let \(G\) be a finite group, and let \(p\) denote the prime implicit in the source's use of mod-\(p\) coefficients. Write \(M=CH_G^*\) for the mod-\(p\) Chow ring referred to in the supplied record. The record does not specify the base field, the construction of this Chow ring, or whether its grading is by codimension or by doubled codimension.

Write \(\mathcal A_p\) for the Steenrod algebra intended by the source. An unstable module here means a graded \(\mathcal A_p\)-module satisfying the instability conditions for the intended algebra and grading. Once those conventions have been fixed, write \(\Sigma^d N\) for the graded suspension, with \((\Sigma^d N)^n=N^{n-d}\) and the corresponding suspended action whenever this is defined in the intended category. Interpret ``contains'' as the existence of an injective, degree-preserving \(\mathcal A_p\)-linear map. Thus the suspension quantity in the source can be expressed as
\[
D(M)=\left\{d\in\Z:\begin{array}{l}
\text{there exists a nonzero unstable module }N,\\
\Sigma^dN\text{ is defined in the intended category, and}\\
\text{there is an injective graded }\mathcal A_p\text{-linear map }
\Sigma^dN\hookrightarrow M
\end{array}\right\}.
\]
The source calls \(d_0(M)\) the topological nilpotence degree, and its literature assessment also calls it a detection bound. No definition of this invariant or its normalization is supplied.

\textbf{Conjecture as recorded.} For a finite group \(G\), the proposed equality is
\[
 d_0(CH_G^*)=\max D(CH_G^*).
\]
This formulation includes the requirement that the indicated maximum exist. The source gives no restriction on \(G\) and states no restriction on the prime \(p\); the intended universal reading is therefore for every finite group and every prime for which the intended construction applies. However, the unspecified Chow-ring setting, Steenrod action, grading, allowable suspension indices, and definition of \(d_0\) must be recovered from the cited material before this becomes a fully determined mathematical statement. In particular, the record does not explicitly say whether only nonnegative suspension indices are allowed.
\subsection{Short English statement}
The supplied conjecture asks whether the topological nilpotence degree of a finite group's mod-\(p\) Chow ring equals the greatest suspension degree of a nonzero unstable Steenrod module that embeds in that ring. The definition of the invariant and several category and grading conventions are missing from the supplied record, so the equality remains conditional on resolving them.
\subsection{Sources}
\begin{itemize}
\item Local UnsolvedMath registry: supplied record with permanent ID \(30002367\), exact source problem code \texttt{catalogue entry}, title ``Topological Nilpotence Degree of Finite-Group Chow Rings,'' in the Oberwolfach Reports Open Problems collection. No individual proposer is supplied; the proposed year is recorded as 2013.
\item Supplied report citation: \emph{Algebraic K-theory and Motivic Cohomology} (2014), Oberwolfach Reports, pp.\ 1861--1913, DOI \texttt{10.4171/owr/2013/32}. \url{https://doi.org/10.4171/owr/2013/32}. The record describes these reports as documenting workshops and their open-problem discussions. The relationship between the citation's 2014 date, the DOI's 2013 component, and the recorded proposed year is not explained.
\item Supplied literature reference, ``Oberwolfach report on Chow detection'': \url{https://ems.press/content/serial-article-files/46461?nt=1}.
\item Supplied literature reference, ``Totaro, Chow ring of a classifying space'': \url{https://arxiv.org/abs/math/9802097}.
\item Supplied literature reference, ``Guillot, Steenrod operations'': \url{https://doi.org/10.1016/j.aim.2004.07.009}.
\item Supplied literature assessment dated \texttt{2026-08-21}: the record labels the problem open and states that the detection bound \(d_0\) is bounded and computed for many finite groups, while the proposed equality remains an open conjecture. This is supplied dated curation derived from the OWR collection, not an independently verified assessment; the record itself requests independent verification of time-sensitive status and cited claims.
\item The registry labels the statement \texttt{corrected\_verified} and explains that notation, typography, and wording were normalized without changing the mathematical content. This supplied editorial label does not establish mathematical correctness or independent source verification.
\end{itemize}
\end{document}
